\chapter{REVIEW OF RELATED LITERATURE AND STUDIES}

\section{Introduction}

This chapter presents the theoretical and empirical foundations relevant to the development of a 2D educational platformer game designed to teach Grade 6 students the concepts of phase changes, heat transfer, and the distinction between physical and chemical changes. The review integrates major learning theories, international research on game-based science education, chemistry-specific digital interventions, and Philippine studies related to educational technology and curriculum alignment.

The chapter is structured into four major sections:

\begin{enumerate}
    \item Theoretical Foundations
    \item International Literature and Studies
    \item Local (Philippine) Literature and Studies
    \item Synthesis and Research Gap
\end{enumerate}

This review establishes the pedagogical legitimacy and scholarly grounding of the developed educational game.

\section{Theoretical Foundations}

\subsection{Conceptual Change Theory}

Conceptual Change Theory (Posner et al., 1982) argues that students enter science classrooms with pre-existing mental models that may conflict with scientifically accepted explanations. Learning, therefore, involves restructuring existing conceptions rather than simply adding new information.

For conceptual change to occur, learners must experience:

\begin{enumerate}
    \item Dissatisfaction with prior conceptions
    \item Intelligibility of the new concept
    \item Plausibility of the new concept
    \item Fruitfulness of the new concept
\end{enumerate}

Research has consistently shown that misconceptions about phase changes are widespread among elementary and middle school students (Driver et al., 1994; Osborne \& Cosgrove, 1983). Common misconceptions include:

\begin{itemize}
    \item Belief that melting alters the substance's identity
    \item Confusion between evaporation and boiling
    \item Misunderstanding of energy transfer during state changes
    \item Belief that physical changes are equivalent to chemical reactions
\end{itemize}

Interactive visual representations have been found to support conceptual restructuring by explicitly modeling particle behavior during phase transitions (Harrison \& Treagust, 2000).

The developed platformer incorporates particle-level manipulation during gameplay. When a player melts an ice block, the lattice structure visually loosens while maintaining particle identity. This directly challenges the misconception that melting changes the substance itself.

Thus, the game design aligns with conceptual change conditions by:

\begin{itemize}
    \item Creating cognitive conflict (blocked path requiring phase change)
    \item Providing intelligible visual models (particle rearrangement)
    \item Demonstrating plausibility (consistent particle identity)
    \item Reinforcing fruitfulness (applying knowledge across levels)
\end{itemize}

\subsection{Constructivism}

Constructivist theory (Piaget, 1970; Bruner, 1966) posits that knowledge is actively constructed through interaction with the environment. Learning is not transmission but transformation.

In science education, constructivist approaches emphasize:

\begin{itemize}
    \item Active engagement
    \item Inquiry-based exploration
    \item Manipulation of representations
    \item Iterative hypothesis testing
\end{itemize}

Digital environments provide rich opportunities for constructivist learning when learners manipulate variables and observe consequences (Jonassen, 1999).

The developed game functions as a constructivist learning space. Players:

\begin{itemize}
    \item Encounter environmental challenges
    \item Experiment with phase transitions
    \item Observe the effects of temperature changes
    \item Apply learned principles in new contexts
\end{itemize}

Rather than presenting static diagrams of molecular motion, the game requires learners to actively reorganize particle systems, aligning with constructivist principles of knowledge construction through experience.

\subsection{Multimedia Learning Theory}

Mayer's Cognitive Theory of Multimedia Learning (Mayer, 2001; Mayer, 2009) proposes that learning is enhanced when information is presented through coordinated visual and verbal channels.

The theory is based on three assumptions:

\begin{enumerate}
    \item Dual-channel processing
    \item Limited working memory
    \item Active processing
\end{enumerate}

Key multimedia principles relevant to the developed game include:

\begin{itemize}
    \item Spatial contiguity principle
    \item Temporal contiguity principle
    \item Coherence principle
    \item Signaling principle
\end{itemize}

Research indicates that animation combined with explanatory text improves understanding of dynamic scientific processes compared to text alone (Mayer \& Moreno, 2002).

In the developed game:

\begin{itemize}
    \item Particle animations are synchronized with labels (e.g., ``Solid $\rightarrow$ Liquid: Melting'')
    \item Visual transformations are immediately tied to terminology
    \item Irrelevant visual elements are minimized during minigames
\end{itemize}

These design features align with multimedia learning principles and strengthen conceptual integration of scientific language with visual representation.

\subsection{Game-Based Learning Theory}

Game-Based Learning (GBL) research emphasizes that games enhance learning when mechanics are directly aligned with instructional objectives (Gee, 2003; Prensky, 2001).

Meta-analyses suggest that digital game-based learning produces moderate positive effects on achievement and motivation in STEM disciplines (Clark et al., 2016; Wouters et al., 2013).

Effective educational games include:

\begin{itemize}
    \item Clear goals
    \item Immediate feedback
    \item Scaffolding
    \item Incremental difficulty
    \item Meaningful player agency
\end{itemize}

The developed platformer integrates:

\begin{itemize}
    \item Goal-driven progression (reach level endpoint)
    \item Mechanic-embedded learning (phase change manipulation required to progress)
    \item Immediate feedback (particle rearrangement animation)
    \item Increasing complexity across levels
\end{itemize}

Unlike superficial quiz-based games, the chemical concept is embedded within the core mechanic of progression.

\subsection{Experiential Learning Theory}

Kolb's Experiential Learning Theory (Kolb, 1984) conceptualizes learning as a cycle involving:

\begin{enumerate}
    \item Concrete experience
    \item Reflective observation
    \item Abstract conceptualization
    \item Active experimentation
\end{enumerate}

Game environments are particularly suited to experiential learning because they allow repeated experimentation with immediate consequences (Squire, 2006).

The platformer structure naturally supports the experiential cycle:

\begin{itemize}
    \item Concrete experience: Encountering a blocked path
    \item Reflection: Observing failure or transformation
    \item Conceptualization: Understanding phase change principles
    \item Experimentation: Applying knowledge in subsequent levels
\end{itemize}

Through repeated cycles, learners internalize abstract chemical principles.

\subsection{Outcome-Based Teaching and Learning}

Outcome-Based Teaching and Learning (OBTL), as articulated by Kalra (2015), emphasizes that educational design should begin with clearly defined learning outcomes, which then determine instructional strategies and assessment methods. Rather than focusing on content delivery alone, OBTL ensures constructive alignment among learning objectives, teaching activities, and assessment practices.

Central principles of OBTL include:

\begin{enumerate}
    \item Clearly articulated measurable learning outcomes
    \item Alignment between outcomes and instructional activities
    \item Assessment strategies that directly measure intended competencies
\end{enumerate}

The MATATAG Curriculum of the Department of Education adopts an outcomes-based orientation, emphasizing competency development, conceptual understanding, and application of knowledge.

The developed educational platformer aligns with OBTL principles by:

\begin{itemize}
    \item Designing gameplay mechanics around specific Grade 6 science competencies
    \item Embedding phase-change manipulation directly into learning objectives
    \item Structuring levels to require demonstration of conceptual understanding
\end{itemize}

Rather than treating the game as supplementary enrichment, the design ensures that progression through levels requires demonstration of intended learning outcomes.

Thus, the game's structure reflects constructive alignment consistent with OBTL principles.

\subsection{Bloom's Taxonomy Framework}

Bloom's Taxonomy (Bloom et al., 1956), later revised by Anderson and Krathwohl (2001), categorizes cognitive processes into hierarchical levels:

\begin{enumerate}
    \item Remember
    \item Understand
    \item Apply
    \item Analyze
    \item Evaluate
    \item Create
\end{enumerate}

The revised taxonomy emphasizes cognitive processes over static knowledge categories.

In the context of Grade 6 chemistry learning, MATATAG competencies require students not only to recall definitions of phase changes but to understand mechanisms, apply concepts, and distinguish physical from chemical changes.

The developed game supports multiple cognitive levels:

\begin{itemize}
    \item \textbf{Remembering:} Identifying terms such as melting, freezing, evaporation
    \item \textbf{Understanding:} Explaining particle movement during state changes
    \item \textbf{Applying:} Using phase changes to solve in-game obstacles
    \item \textbf{Analyzing:} Distinguishing physical changes from chemical changes
\end{itemize}

By embedding conceptual problem-solving into gameplay, the platformer moves beyond rote memorization toward higher-order cognitive engagement.

Table 2.4 presents the alignment between Bloom's cognitive levels, MATATAG competencies, and the developed game mechanics.

\begin{table}[htbp]
\centering
\caption{Alignment of Bloom's Taxonomy, MATATAG Competencies, and Game Mechanics}
\label{tab:blooms_alignment}
\begin{tabularx}{\textwidth}{lXX}
\toprule
\textbf{Bloom's Cognitive Level} & \textbf{MATATAG Competency Focus} & \textbf{Game Implementation} \\
\midrule
Remember & Identify types of phase changes & UI labels and terminology during minigames \\
\addlinespace
Understand & Explain particle behavior during state changes & Animated particle lattice visualization \\
\addlinespace
Apply & Use knowledge of heat transfer in context & Manipulating temperature to overcome obstacles \\
\addlinespace
Analyze & Distinguish physical from chemical changes & Visual preservation of molecular identity \\
\addlinespace
Evaluate (Optional) & Judge appropriate phase change for scenario & Choosing correct state change for progression \\
\bottomrule
\end{tabularx}
\end{table}

\section{Theoretical Framework}

This study is grounded in established learning theories that explain how learners construct, revise, and apply scientific knowledge. The development of the educational platformer game is primarily anchored in Conceptual Change Theory, Constructivist Learning Theory, and Multimedia Learning Theory, with support from Game-Based Learning Theory, Experiential Learning Theory, and Outcome-Based Teaching and Learning principles aligned with Bloom's Taxonomy. These theoretical foundations collectively inform the instructional design, mechanics, and intended learning outcomes of the intervention.

Conceptual Change Theory posits that learners enter the classroom with pre-existing conceptions about scientific phenomena, many of which are misconceptions resistant to change. For conceptual change to occur, learners must experience dissatisfaction with their existing understanding and encounter intelligible, plausible, and fruitful alternative explanations.

In the context of phase changes, learners often hold misconceptions regarding particle motion, energy transfer, and state transitions. The platformer game is designed to confront these misconceptions by allowing learners to manipulate molecular structures directly and observe changes in lattice spacing, particle movement, and energy conditions. By visually and interactively modeling these transitions, the intervention aims to facilitate the restructuring of inaccurate mental models.

Conceptual Change Theory therefore serves as the primary explanatory framework guiding the study's instructional mechanics.

Constructivism asserts that learners actively construct knowledge through interaction with their environment rather than passively receiving information. Learning occurs when individuals engage with tasks, test ideas, and refine understanding through experience.

The game-based intervention reflects constructivist principles by requiring learners to manipulate particle arrangements, experiment with freezing and evaporating mechanics, and solve level-based challenges. Knowledge is not merely presented; instead, learners actively build understanding through interaction and problem-solving.

This theoretical lens justifies the design choice of discovery-based level progression and scaffolded challenges embedded within the platformer environment.

Multimedia Learning Theory suggests that learning improves when information is presented through coordinated visual and verbal channels, provided cognitive processing remains meaningful. Effective multimedia environments integrate text, visuals, and interactive elements to promote deeper understanding.

The platformer incorporates:

\begin{itemize}
    \item Visual representations of particle movement
    \item Explicit textual labels (e.g., ``Solid $\rightarrow$ Liquid: Melting'')
    \item Dynamic changes in structure and spacing
\end{itemize}

These elements are intentionally aligned to reinforce scientific terminology while simultaneously illustrating microscopic processes. The integration of visual simulation and textual labeling supports dual-channel processing and promotes coherent mental model construction.

Game-Based Learning Theory emphasizes motivation, engagement, and structured challenges as mechanisms for promoting learning. The intervention integrates goal-oriented tasks, progressive difficulty, and interactive feedback to sustain learner engagement while targeting specific science competencies.

Experiential Learning Theory proposes that learning occurs through concrete experience, reflective observation, conceptualization, and experimentation. The game enables learners to experience simulated molecular changes, reflect on outcomes, and apply understanding in subsequent levels requiring strategic manipulation of states of matter.

Outcome-Based Teaching and Learning emphasizes alignment between instructional activities and clearly defined learning outcomes. This study aligns game mechanics with competencies specified in the MATATAG Curriculum for Grade 6 Science.

Bloom's Taxonomy further informs the cognitive demands of the game, ensuring that learners move beyond recall toward application and analysis as they manipulate states of matter to solve environmental challenges within the platformer.

Collectively, these theoretical foundations provide the pedagogical justification for the intervention and explain the mechanisms through which conceptual understanding is expected to improve.

\section{Conceptual Framework}

This study adopts an Input--Process--Output (IPO) model to illustrate the relationship between theoretical foundations, the game-based intervention, and the intended learning outcomes.

The conceptual framework synthesizes the theories discussed in the previous section and operationalizes them within the context of an educational platformer game designed to teach phase changes of matter.

\begin{figure}[htbp]
\centering
\includegraphics[width=\textwidth]{ConceptualFrameworkDiagram.png}
\caption{Conceptual Framework of the Study}
\label{fig:conceptual_framework}
\end{figure}

As illustrated in Figure~\ref{fig:conceptual_framework}, the Input component consists of the theoretical and curricular foundations guiding the study. These include Conceptual Change Theory, Constructivism, Multimedia Learning Theory, Game-Based Learning Theory, Experiential Learning Theory, Outcome-Based Teaching and Learning principles, Bloom's Taxonomy, and the MATATAG Curriculum competencies for Grade 6 Science.

These inputs inform the design of the Process component, which is the educational platformer game. The intervention operationalizes theoretical principles through interactive particle manipulation, phase change simulation mechanics, scaffolded level progression, and explicit labeling of state transitions. Learners actively engage with microscopic representations of matter and apply understanding to overcome in-game challenges.

The Output component represents the intended learning outcomes of the study. These include:

\begin{itemize}
    \item Improved conceptual understanding of phase changes
    \item Reduction of misconceptions regarding particle behavior
    \item Enhanced ability to apply knowledge in problem-solving contexts
    \item Alignment with targeted MATATAG Grade 6 Science competencies
\end{itemize}

The arrows connecting the components signify that theoretical and curricular foundations directly shape the instructional design of the intervention, which in turn influences measurable learning outcomes.

This framework guides the development of research instruments and informs the methodology described in Chapter 3.

\section{International Literature and Studies}

\subsection{Digital Game-Based Learning in Science Education}

International studies have documented positive impacts of digital games on science achievement, engagement, and conceptual understanding (Vogel et al., 2006; Wouters et al., 2013).

A meta-analysis by Clark et al. (2016) found that digital game-based learning had a significant positive effect on cognitive outcomes compared to traditional instruction.

Key findings include:

\begin{itemize}
    \item Simulation-based games improve understanding of abstract processes
    \item Interactive environments promote deeper engagement
    \item Embedded feedback enhances conceptual correction
\end{itemize}

To synthesize major international findings, Table 2.1 presents selected studies on digital game-based learning in science education.

\begin{table}[htbp]
\centering
\caption{Summary of Selected International Studies on Game-Based Science Learning}
\label{tab:international_gbl}
\begin{tabularx}{\textwidth}{lllXX}
\toprule
\textbf{Author(s)} & \textbf{Year} & \textbf{Subject Area} & \textbf{Type of Game/Intervention} & \textbf{Key Findings} \\
\midrule
Vogel et al. & 2006 & Various STEM & Computer games \& simulations & Games showed positive effects on cognitive outcomes compared to traditional instruction. \\
\addlinespace
Wouters et al. & 2013 & STEM subjects & Serious games & Serious games improved retention and problem-solving skills, especially with feedback. \\
\addlinespace
Clark et al. & 2016 & Science \& Math & Digital game-based learning & Moderate positive effect size on learning outcomes compared to conventional teaching. \\
\addlinespace
Squire & 2006 & Environmental Science & Immersive simulation game & Context-rich gameplay promoted conceptual understanding and systems thinking. \\
\addlinespace
Gee & 2003 & General learning & Commercial video games & Effective games align mechanics with learning principles such as feedback and scaffolding. \\
\bottomrule
\end{tabularx}
\end{table}

\subsection{Educational Games in Chemistry}

Research specific to chemistry education emphasizes the importance of connecting macroscopic phenomena with submicroscopic particle models (Johnstone, 1991; Treagust et al., 2003).

Studies show that students struggle to transition between:

\begin{itemize}
    \item Observable phenomena
    \item Particle representations
    \item Symbolic notation
\end{itemize}

Interactive simulations that allow manipulation of particle systems have been shown to improve conceptual coherence (Harrison \& Treagust, 2000).

Games that incorporate submicroscopic modeling are particularly effective in addressing misconceptions about phase transitions and intermolecular forces.

The developed game explicitly integrates particle-level manipulation within a platformer mechanic, thereby linking macro-level environmental change (melting wall) with micro-level representation (lattice loosening).

\subsection{Misconceptions in Phase Changes and Heat Transfer}

Extensive research documents misconceptions related to:

\begin{itemize}
    \item Conservation of matter during state change
    \item Role of energy in phase transitions
    \item Distinction between physical and chemical change
\end{itemize}

Students frequently believe that:

\begin{itemize}
    \item Boiling destroys matter
    \item Melting alters chemical composition
    \item Gas particles disappear
\end{itemize}

Interactive digital tools have been recommended to address these misconceptions through dynamic visualization (Driver et al., 1994).

The developed platformer directly addresses these documented conceptual difficulties.

\section{Local (Philippine) Literature and Studies}

\subsection{Educational Technology in Philippine Science Education}

Philippine studies on ICT integration report improved student engagement when interactive tools are incorporated into science instruction (Marcelo, 2012; Caballes, 2016).

However, many implementations involve:

\begin{itemize}
    \item Multimedia presentations
    \item Computer-assisted drills
    \item Video-based instruction
\end{itemize}

Fewer studies examine immersive digital game environments aligned with curriculum competencies.

\subsection{Game-Based Learning in the Philippine Context}

Local research indicates:

\begin{itemize}
    \item Positive student attitudes toward educational games
    \item Increased motivation
    \item Improved short-term retention
\end{itemize}

Yet gaps remain in:

\begin{itemize}
    \item Chemistry-specific educational game research
    \item Particle-level interactive simulations
    \item Curriculum-aligned digital platformer games
\end{itemize}

Table 2.2 summarizes selected Philippine studies related to educational technology and game-based learning.

\begin{table}[htbp]
\centering
\caption{Summary of Selected Philippine Studies on Educational Technology and Game-Based Learning}
\label{tab:philippine_edtech}
\begin{tabularx}{\textwidth}{llllX}
\toprule
\textbf{Author(s)} & \textbf{Year} & \textbf{Level} & \textbf{Subject} & \textbf{Key Findings} \\
\midrule
Marcelo & 2012 & Basic Education & General Science & ICT tools improved engagement but required structured integration for impact. \\
\addlinespace
Caballes & 2016 & Secondary & Science & Technology use increased student participation and conceptual retention. \\
\addlinespace
Various Local Studies & 2015--2020 & Elementary \& Secondary & STEM & Students showed improved motivation; limited chemistry-specific immersive games identified. \\
\bottomrule
\end{tabularx}
\end{table}

\subsection{Alignment with MATATAG Curriculum}

The MATATAG Curriculum emphasizes conceptual understanding and scientific literacy in Grade 6 Science.

Competencies include:

\begin{itemize}
    \item Explaining changes in state of matter
    \item Identifying heat transfer in phase transitions
    \item Distinguishing physical from chemical change
\end{itemize}

The developed game aligns with these competencies by embedding conceptual application within interactive gameplay.

Table 2.3 presents the alignment between MATATAG Grade 6 Science competencies and the core mechanics of the developed game.

\begin{table}[htbp]
\centering
\caption{Alignment of Game Features with MATATAG Grade 6 Science Competencies}
\label{tab:matatag_alignment}
\begin{tabularx}{\textwidth}{XXX}
\toprule
\textbf{MATATAG Competency} & \textbf{Game Mechanic} & \textbf{Intended Learning Outcome} \\
\midrule
Explain changes in state of matter & Particle lattice manipulation minigame & Understand that phase change involves particle rearrangement, not substance change. \\
\addlinespace
Describe the role of heat in phase transitions & Temperature adjustment interaction & Recognize heat transfer as the driver of melting, freezing, evaporation. \\
\addlinespace
Distinguish physical and chemical changes & Visual preservation of particle identity & Identify that chemical composition remains constant during phase change. \\
\addlinespace
Apply understanding of matter in problem-solving & Platform obstacles requiring phase manipulation & Transfer conceptual knowledge to new scenarios. \\
\bottomrule
\end{tabularx}
\end{table}

\section{Synthesis of Literature}

The reviewed literature reveals several consistent findings:

\begin{enumerate}
    \item Students possess persistent misconceptions regarding phase changes.
    \item Particle-level visualization improves conceptual understanding.
    \item Game-based learning enhances engagement and cognitive outcomes.
    \item Alignment between mechanics and learning objectives is critical.
    \item Philippine research lacks immersive chemistry-focused platformer interventions.
\end{enumerate}

Collectively, these findings support the pedagogical rationale for developing an interactive educational platformer grounded in conceptual change and constructivist principles.

\section{Research Gap}

Despite strong evidence supporting digital game-based learning in science education, the following gaps remain:

\begin{itemize}
    \item Limited research integrating submicroscopic particle modeling into platformer mechanics.
    \item Few curriculum-aligned chemistry educational games designed for Grade 6 learners in the Philippine context.
    \item Insufficient research examining the impact of interactive lattice manipulation on correcting misconceptions regarding physical versus chemical change.
\end{itemize}

This study addresses these gaps by developing and evaluating a curriculum-aligned educational platformer embedding conceptual change mechanisms within gameplay.
